\section{PROCEDIMENTOS}

\emph{A. Fotodiodo}

Realiza-se a montagem do circuito abaixo e levanta-se a curva característica IDxVD, variando-se a tensão na fonte de 0 V a 5 V em passos de 200 mV.

% Imagem Circ 1 - font -> resistência -> led

\emph{B. Fototransistor}

Monta-se o circuito abaixo cuja a resistência é controlável pela potência luminosa incidente sobre ele numa dada faixa do espectro de frequências. Em seguida, levanta-se a característica ICxVce do fototransistor TIL78 nas seguintes condições: Sob luz incidente, com iluminação extra de uma lâmpada controlada por um varivolt em 60 mV, mesma iluminação extra porém com o varivolt em 80 mV. Para levantar a curva nas condições citadas acima, deve-se variar a tensão na fonte de 0 a 6 V em passos de 0,5 V, determinando a corrente no dispositivo e a tensão entre seus terminais em cada caso.

% Imagem Circ 2 - font -> resistência -> fotoresistor 

\emph{C. Fotocélulas}

Com o uso de fotocélulas, mede-se a tensão de circuito aberto e a corrente de curto-circuito em diferentes condições de iluminação incidente na fotocélula.

\emph{D. Acoplamento óptico}

Monta-se o circuito abaixo de acoplamento óptico e, colocando-se uma tensão V1 na fonte, verifica-se o comportamento captado pelo fototransistor através do osciloscópio.

% Imagem Circ 3 - acoplamento óptico (transmissão de sinais entre sistemas eletrônicos por meio de luz)

\vspace{.25cm}